\documentclass[11pt]{article}
\usepackage[T1]{fontenc}
\usepackage{lmodern}
\usepackage[margin=1in]{geometry}
\usepackage{amsmath,amssymb,amsthm}
\usepackage{microtype}
\usepackage[colorlinks=true,linkcolor=blue,citecolor=blue,urlcolor=blue]{hyperref}
\newtheorem{theorem}{Theorem}
\newtheorem{lemma}[theorem]{Lemma}
\newtheorem{proposition}[theorem]{Proposition}
\newtheorem{remark}[theorem]{Remark}
\newcommand{\R}{\mathbb R}
\newcommand{\E}{\mathbb E}
\newcommand{\Prob}{\mathbb P}
\newcommand{\HS}{\mathrm{HS}}
\newcommand{\off}{\operatorname{off}}
\newcommand{\diag}{\operatorname{diag}}
\newcommand{\Tr}{\operatorname{Tr}}
\newcommand{\rank}{\operatorname{rank}}
\newcommand{\ind}{\mathbf 1}
\title{Least singular values of independent random tensor columns}
\author{Kaiwen Jin\\
\small\href{mailto:yui.ui.siki@gmail.com}{yui.ui.siki@gmail.com}}
\date{October 13, 2026}
\begin{document}
\maketitle

\begin{abstract}
Let \(A\) have \(m=\lfloor(1-\varepsilon)n^d\rfloor\) independent
columns, each the tensor product of \(d\) random vectors in \(\R^n\).
All scalar coordinates are independent, centered, of variance one,
and have subgaussian norm at most \(K\). We prove that, for every fixed
\(\varepsilon\in(0,1/2]\) and every integer sequence
\(d=o(\sqrt{n/\log n})\), the least singular value is at least
\(c_K\varepsilon n^{d/2}\) with probability tending to one.
The constant depends only on \(K\); neither identical distributions
nor symmetry of the coordinates is required. The proof conditions on
all but the last tensor factor, estimates quadratic Gram moments of
projected column blocks, and lifts the block estimate by averaging
over subsets of columns.
\end{abstract}

\section{Introduction}

Independent rank-one tensor features arise when coefficients are recovered
from a linear combination of simple tensors. If these features are the
columns of a matrix \(A\), a lower bound for its least singular value controls
the amplification of an error in the observed linear combination.
Tensor columns are isotropic under the assumptions below, but their scalar
entries are dependent. The tensor order may also grow with the factor
dimension.

For a real random variable \(\xi\), use the convention
\[
 \|\xi\|_{\psi_2}
 =\inf\{s>0:\E\exp(\xi^2/s^2)\le2\}.
\]
Let \(x_{i\ell}\in\R^n\), \(1\le i\le m\), \(1\le\ell\le d\).
Assume that all coordinates of all these vectors are mutually independent,
with mean zero, variance one, and subgaussian norm at most \(K\).
Define
\[
 X_i=x_{i1}\otimes\cdots\otimes x_{id},\qquad
 A=[X_1,\ldots,X_m]\in\R^{N\times m},\qquad N=n^d.
\]
The coordinates can have different laws, and the laws can depend on \(n\).
For a matrix with \(m\ge1\) columns, write
\[
 s_{\min}(A)=\inf\{\|Av\|:v\in\R^m,\ \|v\|=1\}.
\]

Vershynin's tensor concentration theorem has constants independent of
the order and yields quantitative linear independence
\cite[Theorem 1.4 and Corollary 6.2]{Vershynin}.
In the regime
\(C_Kd^2\log(n)/n\le\varepsilon\le1/2\), that corollary gives
\(s_{\min}(A)\ge\sqrt{\varepsilon}/2\) with probability at least
\(1-2\exp(-c_K\varepsilon n/d)\), when
\(m\le(1-\varepsilon)n^d\).
Section 1.5 of the same paper asks for an improvement to the ambient
dimension scale. The following theorem supplies that scale with a constant
independent of the tensor order.

\begin{samepage}
\begin{theorem}\label{thm:main}
For every \(K>0\) there is \(c_K>0\) with the following property.
Fix any \(\varepsilon\in(0,1/2]\), and let \(d=d(n)\ge1\) be an integer
sequence such that
\[
 \frac{d^2\log n}{n}\longrightarrow0.
\]
Put \(N=n^d\) and \(m=\lfloor(1-\varepsilon)N\rfloor\).
For every sequence of coordinate laws satisfying the independence,
centering, variance and subgaussian assumptions above,
\[
 \Prob\{s_{\min}(A)\ge c_K\varepsilon\sqrt N\}
 \longrightarrow1\qquad(n\longrightarrow\infty).
\]
The constant \(c_K\) is independent of \(n,d,\varepsilon\) and the
coordinate laws. Probability convergence is asserted for each fixed
\(\varepsilon\).
\end{theorem}
\end{samepage}

This is the full target of Problem 562 in the AIM collection
\cite{AIM}, at the revision specified in that reference.

An estimate of the distance of each column from the span of the others
does not by itself give this scale: selecting a largest coordinate of a
unit coefficient vector can lose a factor \(\sqrt m\).
We instead project a small block of \(k\) columns away from all remaining
columns. Its Gram matrix is well conditioned at the projection-rank
scale. Once this holds with high probability for each fixed block,
a deterministic subset-averaging argument gives the conclusion for all
coefficient vectors without a union bound over all blocks.

To estimate a projected block, we condition on its first \(d-1\) factors.
The remaining Gram fluctuation is a quadratic polynomial in independent
scalars. A scalar convex comparison reduces its off-diagonal part to a
Gaussian quadratic polynomial with constants depending only on \(K\).
We then prove a Gaussian Gram moment bound using Hermite creation and
annihilation operators. This use of low-degree operator norms follows the
method of Bandeira, Kunisky, Nizi\'c-Nikolac, Pesenti and Wang
\cite[Section 2]{Fock}; the coefficient identities and estimates needed
here are proved in Section~\ref{sec:moments}. Tensor concentration
\cite{Vershynin} and matrix Chernoff \cite{Tropp} control the conditioned
coefficients. These are the two concentration estimates imported into
the proof.

All vector norms are Euclidean and all matrix norms are induced operator
norms. The Hilbert--Schmidt norm is denoted by \(\|\cdot\|_{\HS}\),
and \(\Tr\) is the unnormalized trace. For a square matrix \(H\),
\(\off(H)=H-\diag(H)\). All logarithms are natural.

\section{Quadratic Gram moments}\label{sec:moments}

\subsection{Scalar comparison}

We first record an elementary scalar form of convex domination.
For broader vector comparison results, see van Handel \cite{vanHandel}.

\begin{lemma}\label{lem:convex}
If \(\E\xi=0\) and \(\|\xi\|_{\psi_2}\le K\), and \(g\) is standard
normal, then
\[
 \E f(\xi)\le\E f(16K g)
\]
for every finite convex function \(f:\R\to\R\) of polynomial growth.
\end{lemma}

\begin{proof}
Put \(Y=\xi/(2K)\). Then \(\E e^{Y^2}\le2\), so
\(\Prob\{|Y|>u\}\le2e^{-u^2}\). The inequality
\(Y^2\le e^{Y^2}-1\) gives \(\E Y^2\le1\), and hence
\(\E|Y|\le1\).
For \(0\le z\le2\),
\[
 \E(Y-z)_+\le1
 <\frac8{\sqrt{2\pi}}-2
 \le\E(8g-z)_+.
\]
For \(z\ge2\), integration of the upper tail gives
\[
 \E(Y-z)_+\le2\int_z^\infty e^{-u^2}\,du
 \le\frac{e^{-z^2}}{z}\le e^{-z^2}.
\]
On the event \(g\ge(z+1)/8\), the positive part \((8g-z)_+\) is
at least one. Integrating the Gaussian density over the interval
\([(z+1)/8,(z+9)/8]\) therefore gives
\[
 \E(8g-z)_+
 \ge\Prob\{g\ge(z+1)/8\}
 \ge\frac1{\sqrt{2\pi}}\exp\!\left(-\frac{(z+9)^2}{128}\right)
 \ge e^{-z^2}.
\]
The last inequality holds because
\(z^2-(z+9)^2/128\) is increasing for \(z\ge2\), and its value
at \(2\) is greater than \(3>\log\sqrt{2\pi}\).
The same bounds hold for \(-Y\). Since \(Y\) is centered,
\(\E(Y-z)_+=-z+\E(-Y+z)_+\); this proves the comparison of
positive parts for negative \(z\) as well.

A finite convex piecewise affine function is an affine function plus
a finite nonnegative linear combination of functions \((x-z)_+\).
The comparison just proved and centering give its expectation inequality.
For a general finite convex \(f\), choose one supporting line at each
rational point and include a fixed supporting line in every finite
maximum. These maxima increase to \(f\): supporting slopes are locally
bounded by secant slopes, so rational points approaching any point
recover its value. Subtracting the fixed line leaves a nonnegative
increasing sequence. Monotone convergence proves the expectation
comparison for \(Y\) and \(8g\). Applying this to \(f(2K\,\cdot)\)
proves the claim. The polynomial-growth assumption ensures that all
expectations are finite.
\end{proof}

\begin{lemma}\label{lem:quadratic}
Let \(\xi_1,\ldots,\xi_s\) be independent centered scalars with
\(\|\xi_\alpha\|_{\psi_2}\le K\). Let \(H_{\alpha\lambda}\)
be real symmetric matrices of the same size, with
\(H_{\alpha\lambda}=H_{\lambda\alpha}\) and \(H_{\alpha\alpha}=0\).
For independent standard normals \(g_\alpha\), and any \(q\ge1\),
\[
 \left(\E\left\|\sum_{\alpha,\lambda}
 H_{\alpha\lambda}\xi_\alpha\xi_\lambda\right\|^q\right)^{1/q}
 \le2(16K)^2
 \left(\E\left\|\sum_{\alpha,\lambda}
 H_{\alpha\lambda}g_\alpha g_\lambda\right\|^q\right)^{1/q}.
\]
\end{lemma}

\begin{proof}
Write \(Q(x)=\sum_{\alpha,\lambda}H_{\alpha\lambda}
x_\alpha x_\lambda\). Include each scalar index in a random subset
\(S\) independently with probability \(1/2\), and let
\[
 Q_S(x)=\sum_{\alpha\in S,\ \lambda\notin S}
 H_{\alpha\lambda}x_\alpha x_\lambda.
\]
Then \(Q(x)=4\E_SQ_S(x)\), so convexity of the norm power yields
\[
 \|Q(\xi)\|_{L^q}
 \le4\big(\E_S\E\|Q_S(\xi)\|^q\big)^{1/q}.
\]
For fixed \(S\), the expression \(Q_S\) is affine in each separate
coordinate. Conditional on the others, its norm to the power \(q\)
is a convex function of that coordinate. Successive application of
Lemma~\ref{lem:convex}, followed by degree-two homogeneity, gives
\[
 \|Q_S(\xi)\|_{L^q}\le(16K)^2\|Q_S(g)\|_{L^q}.
\]
If \(g^S\) changes the signs of the coordinates indexed by \(S\),
symmetry of the coefficient array implies
\(Q_S(g)=(Q(g)-Q(g^S))/4\).
Since \(g^S\) has the same law as \(g\), the \(L^q\) triangle
inequality gives
\(\|Q_S(g)\|_{L^q}\le\|Q(g)\|_{L^q}/2\).
Combining the bounds proves the lemma.
\end{proof}

\subsection{Gaussian column matrices}

The next estimate keeps both covariance scales separate. This is what
makes a small projected block useful.

\begin{lemma}\label{lem:gaussian}
Let \(U_i:\R^n\to\R^M\), \(1\le i\le k\), be deterministic real
matrices, and let \(g_i\) be independent standard Gaussian vectors in
\(\R^n\). Put
\[
 W=[U_1g_1,\ldots,U_kg_k],\quad
 D=\E W^TW=\diag(\|U_i\|_{\HS}^2),\quad
 L=\E WW^T=\sum_iU_iU_i^T,
\]
and write \(a_0=\|D\|\), \(b_0=\|L\|\), and
\(t=\max_i\|U_i\|^2\).
For every integer \(p\ge1\), set
\[
 \mathcal M_p=b_0+2(2p+1)t
       +2\sqrt{(a_0+(2p+1)t)(b_0+(2p+1)t)}.
\]
Then
\begin{equation}\label{eq:gaussian-moment}
 \big(\E\|\off(W^TW)\|^{2p}\big)^{1/(2p)}
 \le2k^{1/(2p)}\mathcal M_p.
\end{equation}
\end{lemma}

\begin{proof}
Denote the columns of \(U_i\) by \(u_{ia}\), index scalar
coefficients by \(\alpha=(i,a)\), and set
\(C_\alpha=u_{ia}e_i^T\in\R^{M\times k}\).
Thus \(W=\sum_\alpha g_\alpha C_\alpha\).
We need the two self-adjoint block matrices
\begin{equation}\label{eq:blocks}
 R_{\alpha\lambda}=C_\lambda^TC_\alpha,\qquad
 S_{\alpha\lambda}=C_\lambda C_\alpha^T.
\end{equation}
They satisfy
\begin{equation}\label{eq:block-norms}
 \|R\|\le t,\qquad \|S\|\le t.
\end{equation}
Indeed, for an input with coordinates \(f_{(h,b);i}\), the output
of \(R\) is
\[
 (Rf)_{(i,a);h}
   =\sum_b\langle u_{hb},u_{ia}\rangle f_{(h,b);i}.
\]
For each ordered pair \((i,h)\), this is the map \(U_i^TU_h\)
from the input block labelled \((h,i)\) to the output block
labelled \((i,h)\). Both block families are disjoint, and each
map has norm at most \(t\). Summing squared norms proves the first
bound. For \(S\), coefficients with distinct column indices vanish.
For each \(i\), the block matrix \([u_{ib}u_{ia}^T]_{a,b}\) is
\[
 (I_n\otimes U_i)F(I_n\otimes U_i^T),
\]
where \(F\) interchanges the two \(\R^n\) factors. This gives the
second bound.

We use the Gaussian \(L^2\) inner product on polynomials.
The Hermite polynomials defined by
\(H_0=1\), \(H_1(x)=x\), and
\(H_{\ell+1}=xH_\ell-\ell H_{\ell-1}\) satisfy
\(H_\ell'=\ell H_{\ell-1}\).
For the standard Gaussian density \(\phi\),
\(\phi^{(\ell)}=(-1)^\ell\phi H_\ell\).
Integration by parts \(\ell\) times shows that
\(\E H_\ell(g)H_j(g)=\E H_j^{(\ell)}(g)\).
This is zero if \(\ell>j\), and equals \(\ell!\) if \(\ell=j\);
the other unequal case follows by symmetry.
The boundary terms vanish because they are polynomials times \(\phi\).
The normalized polynomials \(h_\ell=H_\ell/\sqrt{\ell!}\) therefore
form an orthonormal basis of the polynomial space. Their recurrence is
\[
 xh_\ell(x)=\sqrt{\ell+1}\,h_{\ell+1}(x)
                    +\sqrt{\ell}\,h_{\ell-1}(x).
\]
For \(\eta\in\mathbb N_0^{kn}\), the products
\(h_\eta=\prod_\alpha h_{\eta_\alpha}\) form the corresponding
multivariate orthonormal basis. Define creation and annihilation on
this basis by
\[
 c_\alpha h_\eta=\sqrt{\eta_\alpha+1}\,h_{\eta+e_\alpha},
 \qquad
 c_\alpha^*h_\eta=\sqrt{\eta_\alpha}\,h_{\eta-e_\alpha}.
\]
A term with \(\eta_\alpha=0\) in the second formula is zero.
Multiplication by \(g_\alpha\) is \(c_\alpha+c_\alpha^*\), and
\[
 c_\lambda^*c_\alpha-c_\alpha c_\lambda^*
   =\ind_{\alpha=\lambda}I.
\]
Let \(\mathcal P_{\le q}\) be the polynomials of total degree at
most \(q\), and let \(\Lambda f=(c_\alpha^*f)_\alpha\).
On \(\mathcal P_{\le q}\), also for vector-valued polynomials,
\begin{equation}\label{eq:annihilation}
 \|\Lambda\|\le\sqrt q,
\end{equation}
since \(\sum_\alpha c_\alpha c_\alpha^*h_\eta=|\eta|h_\eta\).

Define
\(\mathcal U=\sum_\alpha C_\alpha c_\alpha\) and
\(\mathcal V=\sum_\alpha C_\alpha c_\alpha^*\).
These map \(k\)-component polynomials to \(M\)-component polynomials,
and multiplication by \(W\) is \(\mathcal U+\mathcal V\).
The commutator and \eqref{eq:blocks} give the identities
\begin{equation}\label{eq:operator-identities}
 \mathcal U^*\mathcal U=D+\Lambda^*R\Lambda,\qquad
 \mathcal V\mathcal V^*=L+\Lambda^*S\Lambda.
\end{equation}
Consequently, on degree at most \(q\),
\[
 \|\mathcal U\|\le\sqrt{a_0+qt},\qquad
 \|\mathcal V^*\|\le\sqrt{b_0+qt}.
\]
On degree at most \(q\ge1\), their degree-lowering adjoints satisfy
\[
 \|\mathcal U^*\|\le\sqrt{a_0+(q-1)t},\qquad
 \|\mathcal V\|\le\sqrt{b_0+(q-1)t}.
\]
For these last bounds, restrict the raising maps to degree at most
\(q-1\) and take adjoints; homogeneous degrees are orthogonal.
Both lowering maps vanish on constants.
All norms here are restrictions between finite-dimensional
polynomial spaces. No boundedness of these operators on the full
Gaussian \(L^2\) space is asserted.

Set \(T=W^TW-D\), viewed also as a multiplication operator.
By \eqref{eq:operator-identities},
\[
 T=\Lambda^*R\Lambda+
       \mathcal U^*\mathcal V+
       \mathcal V^*\mathcal U+
       \mathcal V^*\mathcal V.
\]
On input degree at most \(q\), the first term has norm at most
\(qt\). Each middle term has norm at most
\(\sqrt{(a_0+(q+1)t)(b_0+(q+1)t)}\), and the last has norm at
most \(b_0+(q+1)t\). These bounds also hold when \(q=0\), by the
vanishing of the lowering maps on constants.
Thus \(T\) raises degree by at most two and its restricted norm is
at most
\[
 b_0+(2q+1)t+
 2\sqrt{(a_0+(q+1)t)(b_0+(q+1)t)}.
\]
In forming \(T^p(e_i\otimes1)\), the input degrees are
\(0,2,\ldots,2p-2\); the last display is bounded by \(\mathcal M_p\)
at each of these steps. Therefore
\[
 \E\Tr(T^{2p})
 =\sum_{i=1}^k\|T^p(e_i\otimes1)\|_{L^2}^2
 \le k\mathcal M_p^{2p}.
\]
Pointwise symmetry gives \(\|T\|^{2p}\le\Tr(T^{2p})\).
Finally,
\(\off(W^TW)=T-\diag(T)\) and
\(\|\diag(T)\|\le\|T\|\). This proves
\eqref{eq:gaussian-moment}.
\end{proof}

\section{Tensor blocks and the full matrix}\label{sec:tensors}

We now apply the moment bound in a way that keeps its constants
independent of \(d\). We use the following form of
\cite[Theorem 1.4]{Vershynin}.
If \(Y\) is a tensor of order \(q\ge1\) with independent, centered,
unit-variance, \(K\)-subgaussian coordinates in its factors, and \(B\)
is a deterministic linear map from \(\R^{n^q}\) to a Euclidean
space, then
\begin{equation}\label{eq:tensor-concentration}
 \Prob\{|\|BY\|-\|B\|_{\HS}|\ge z\}
 \le2\exp\!\left(-\frac{a_K z^2}
                       {q n^{q-1}\|B\|^2}\right),
 \qquad 0\le z\le2\|B\|_{\HS}.
\end{equation}
The constant \(a_K>0\) depends only on \(K\).
Zero maps can be treated deterministically.
We also use the upper tail in
\cite[Theorem 5.1.1]{Tropp}: for independent positive semidefinite
\(r_0\times r_0\) matrices \(Y_i\le L_0I\), with
\(\mu=\|\sum_i\E Y_i\|>0\),
\begin{equation}\label{eq:cher}
 \Prob\!\left\{\left\|\sum_iY_i\right\|
                    \ge(1+\theta)\mu\right\}
 \le r_0\left[\frac{e^\theta}
                   {(1+\theta)^{1+\theta}}\right]^{\mu/L_0},
 \qquad \theta\ge0.
\end{equation}

\subsection{Projected blocks}

\begin{proposition}\label{prop:block}
Under the assumptions of Theorem~\ref{thm:main}, there is
\(\delta_K\in(0,1/4)\), depending only on \(K\), as follows.
Put \(h=N-m\) and \(k=\lfloor\delta_K h\rfloor\).
For each fixed \(k\)-subset \(J\) of columns, let \(P_J\) project
orthogonally off the span of the columns outside \(J\).
For sufficiently large \(n\), \(1\le k\le m\), and
\[
 \Prob\{s_{\min}(P_JA_J)<\sqrt{\varepsilon N/8}\}
 \le q_n,\qquad q_n\longrightarrow0.
\]
The bound \(q_n\) can be chosen independently of \(J\) and the
coordinate laws.
\end{proposition}

\begin{proof}
For now let \(\delta\in(0,1/4)\) be a constant to be chosen using
only \(K\), and set \(k=\lfloor\delta h\rfloor\).
For fixed \(\varepsilon\) and sufficiently large \(n\),
\[
 k\ge\delta h/2,\qquad 1\le k\le m.
\]
Fix \(J\) and condition on all columns outside \(J\).
Write \(P=P_J\), \(r=\rank P\), and
\[
 \beta=r/N,\qquad\gamma=k/N.
\]
There are \(m-k\) outside columns, even if some are dependent; hence
\begin{equation}\label{eq:rank}
 r\ge N-(m-k)=h+k,\qquad
 \beta\ge\varepsilon,\qquad\gamma\le\delta\beta.
\end{equation}
The projection is measurable without any rank assumption on those
columns. If their matrix is \(C\), it is the pointwise limit
\[
 P=\lim_{\lambda\downarrow0}
 \big[I-C(C^TC+\lambda I)^{-1}C^T\big].
\]
The singular-value decomposition verifies the limit. In particular,
\(P\) is independent of all coordinates of columns in \(J\), so
conditioning means integration over the outside product law.

Let \(R_0=n^{d-1}\) and
\(b_i=x_{i1}\otimes\cdots\otimes x_{i,d-1}\in\R^{R_0}\).
For \(d=1\), set \(R_0=1\) and \(b_i=1\).
After an orthogonal reordering of tensor coordinates, there are
isometric embeddings \(E_a:\R^{R_0}\to\R^N\) into \(n\)
mutually orthogonal blocks such that
\[
 X_i=\sum_{a=1}^n \xi_{ia}E_ab_i,
\]
where the \(\xi_{ia}\) are the coordinates of the last factor.
Condition next on all \(b_i\), \(i\in J\), and define
\begin{equation}\label{eq:Ui}
 U_i=\frac{[PE_1b_i,\ldots,PE_nb_i]}{\sqrt N},
 \qquad W=\frac{PA_J}{\sqrt N}.
\end{equation}

We first establish, with high probability over the \(b_i\),
\begin{equation}\label{eq:typical}
 \max_{i\in J}\|b_i\|^2\le2R_0,\qquad
 \left\|\sum_{i\in J}b_ib_i^T\right\|\le2k,\qquad
 \max_{i\in J}\|U_i\|_{\HS}^2\le9\beta/4.
\end{equation}
For \(d\ge2\), apply \eqref{eq:tensor-concentration} to the
identity on order \(d-1\) tensors, with
\(z=(\sqrt2-1)\sqrt{R_0}\), for the first bound.
For the second, truncate separately in each column:
\[
 Y_i=b_ib_i^T\ind_{\{\|b_i\|^2\le2R_0\}}.
\]
These matrices are independent, \(0\le Y_i\le2R_0I\), and
\(\E Y_i\le I\).
To see the last assertion, note that \(\E b_ib_i^T=I\):
for tensor multi-indices, the covariance factors into scalar
covariances and equals one on the diagonal and zero otherwise.
This uses independence and unit variance, with no equality of laws.
Let \(\mu=\|\sum_i\E Y_i\|\le k\).
If \(\mu=0\), positivity and a zero expected trace give
\(\sum_iY_i=0\) almost surely.
Otherwise take \(\theta=2k/\mu-1\ge1\) in \eqref{eq:cher}.
The logarithm of the bracketed probability factor is
\[
 \frac{2k-\mu-2k\log(2k/\mu)}{2R_0}.
\]
This is increasing in \(\mu\) for \(0<\mu\le k\), and its value
at \(k\) is \(-(2\log2-1)k/(2R_0)\). Therefore
\begin{equation}\label{eq:covariance-tail}
 \Prob\!\left\{\left\|\sum_iY_i\right\|>2k\right\}
 \le R_0\exp\!\left(-\frac{(2\log2-1)k}{2R_0}\right).
\end{equation}
On the first event in \eqref{eq:typical}, \(\sum_iY_i=\sum_ib_ib_i^T\).

For the third bound in \eqref{eq:typical}, put
\(T_P=\sum_aE_a^TPE_a\). Then \(0\le T_P\le nI\) and
\(\Tr T_P=r\). Apply \eqref{eq:tensor-concentration} to
\(T_P^{1/2}\) and \(b_i\), with \(z=\sqrt r/2\).
Since
\(\|U_i\|_{\HS}^2=b_i^TT_Pb_i/N\), this gives the desired
upper bound. All these applications are uniform in the fixed
projection \(P\).
For \(d=1\), the three bounds hold deterministically:
\(\|b_i\|^2=1\), \(\sum_ib_ib_i^T=k\), and
\(\|U_i\|_{\HS}^2=\beta\).

On \eqref{eq:typical}, the parameters of Lemma~\ref{lem:gaussian}
for these \(U_i\) satisfy
\begin{equation}\label{eq:parameters}
 a_0\le9\beta/4,\qquad b_0\le2\gamma,\qquad t\le2/n.
\end{equation}
For the middle inequality,
\[
 \sum_iU_iU_i^T
 =\frac{P(I_n\otimes\sum_ib_ib_i^T)P}{N}.
\]
For the last, the unprojected columns \(E_ab_i\) are orthogonal,
so \(\|U_i\|\le\|b_i\|/\sqrt N\).

The off-diagonal Gram matrix in \eqref{eq:Ui} is a quadratic
polynomial to which Lemma~\ref{lem:quadratic} applies.
Explicitly, set \(C_{(i,a)}=PE_ab_i e_i^T/\sqrt N\).
For \(\alpha=(i,a)\), \(\lambda=(j,b)\), take
\[
 H_{\alpha\lambda}=
 \begin{cases}
 (C_\alpha^TC_\lambda+C_\lambda^TC_\alpha)/2,&i\ne j,\\
 0,&i=j.
 \end{cases}
\]
These coefficients are symmetric as matrices and in their two
indices, and
\(\off(W^TW)=\sum_{\alpha,\lambda}
H_{\alpha\lambda}\xi_\alpha\xi_\lambda\).
The last-factor coordinates retain their independent centered
laws after the conditioning. Lemmas~\ref{lem:quadratic}
and \ref{lem:gaussian} therefore give
\[
 \|\off(W^TW)\|_{L^{2p}}
 \le4(16K)^2k^{1/(2p)}\mathcal M_p,
\]
where this moment is conditional on \(P,(b_i)\).
Choose \(p=\lceil\log N\rceil\), and put
\(\rho=2(2p+1)/n\) and \(B_K=4(16K)^2\sqrt e\).
Since \(k\le N\), \eqref{eq:parameters} bounds the last display by
\begin{equation}\label{eq:conditional-moment}
 B_K\left[2\gamma+2\rho+
 2\sqrt{(9\beta/4+\rho)(2\gamma+\rho)}\right].
\end{equation}
Choose \(\delta=\delta_K\in(0,1/4)\) so that
\begin{equation}\label{eq:delta}
 B_K[2\delta_K+3\sqrt{2\delta_K}]<1/(16e).
\end{equation}
This choice uses only \(K\).
Dividing \eqref{eq:conditional-moment} by \(\beta\),
using \eqref{eq:rank}, and setting \(\eta=\rho/\beta\)
gives the upper bound
\[
 B_K\left[2\delta_K+2\eta+
 2\sqrt{(9/4+\eta)(2\delta_K+\eta)}\right].
\]
At \(\eta=0\), this is less than \(1/(16e)\).
By continuity it is at most \(1/(8e)\) for
\(0\le\eta\le\eta_K\), for some \(\eta_K>0\) depending only on \(K\).
Since \(\rho/\varepsilon\to0\) and \(\beta\ge\varepsilon\),
this holds for all sufficiently large \(n\). Thus, on
\eqref{eq:typical}, Markov's inequality gives
\begin{equation}\label{eq:off-tail}
 \Prob\{\|\off(W^TW)\|>\beta/8\mid P,(b_i)\}
 \le e^{-2p}\le N^{-2}.
\end{equation}

For the actual diagonal entries, condition only on \(P\).
Apply \eqref{eq:tensor-concentration} at order \(d\), with
\(B=P\), \(\|P\|_{\HS}=\sqrt r\), and \(z=\sqrt r/2\).
A union bound gives
\begin{equation}\label{eq:diag}
 \min_{i\in J}\frac{\|PX_i\|^2}{N}\ge\beta/4
\end{equation}
except with probability at most \(2k\exp(-a_K\beta n/d)\),
after decreasing \(a_K\) by an absolute factor.
There is no need for independence between \eqref{eq:diag}
and the event in \eqref{eq:off-tail}.
Their intersection implies
\(\lambda_{\min}(W^TW)\ge\beta/8\), and hence
\[
 s_{\min}(PA_J)\ge\sqrt{\beta N/8}
                  \ge\sqrt{\varepsilon N/8}.
\]

Combining the three tensor concentration bounds, the covariance
tail and \eqref{eq:off-tail}, the failure probability is at most
\begin{equation}\label{eq:qn}
 q_n=6N\exp(-a_K\varepsilon n/d)
       +R_0\exp(-c_0\delta_K\varepsilon n)+N^{-2},
 \qquad c_0=(2\log2-1)/4,
\end{equation}
for sufficiently large \(n\), with a further decrease of \(a_K\)
depending only on \(K\).
Here \(k/R_0\ge\delta_K\varepsilon n/2\).
The first term tends to zero because
\[
 \frac{d\log n}{\varepsilon n/d}
       =\frac{d^2\log n}{\varepsilon n}\longrightarrow0,
\]
the second because \(\log R_0=(d-1)\log n=o(n)\),
and the third because \(N\to\infty\).
Also \(\rho/\varepsilon=O(d\log n/(\varepsilon n))\to0\),
as used above.
For \(d=1\), the auxiliary events in \eqref{eq:typical} are
deterministic, so the same upper bound remains valid.
Every bound is uniform in \(J\), the conditioned projection,
and the permitted laws.
\end{proof}

\subsection{Subset averaging}

We conclude with a deterministic implication for any matrix.
For a fixed unit vector \(v\in\R^m\), let \(J\) be uniform among
the \(k\)-subsets of \(\{1,\ldots,m\}\), and write
\[
 Z=\|v_J\|^2,\qquad t_0=k/m.
\]
Then \(0\le Z\le1\) and \(\E_JZ=t_0\). It follows that
\begin{equation}\label{eq:subset-mass}
 \Prob_J\{Z\ge t_0/2\}
 \ge\frac{t_0}{2-t_0}\ge t_0/2,
\end{equation}
by bounding \(\E_JZ\) above by
\(t_0/2+(1-t_0/2)\Prob_J\{Z\ge t_0/2\}\).

\begin{proof}[Proof of Theorem~\ref{thm:main}]
Take \(\delta_K\) from Proposition~\ref{prop:block}, and set
\[
 c_K=\sqrt{\delta_K}/8,\qquad
 k=\lfloor\delta_K(N-m)\rfloor.
\]
For sufficiently large \(n\),
\[
 t_0=k/m\ge\delta_K\varepsilon/2.
\]
For any realization of \(A\) with
\(s_{\min}(A)<c_K\varepsilon\sqrt N\), choose a unit vector \(v\)
with \(\|Av\|<c_K\varepsilon\sqrt N\).
If \(Z\ge t_0/2\) and the projected block for \(J\) obeyed
\(s_{\min}(P_JA_J)\ge\sqrt{\varepsilon N/8}\), then
\[
 \|Av\|\ge\|P_JAv\|=\|P_JA_Jv_J\|
 \ge\sqrt{\varepsilon N/8}\sqrt{t_0/2}
 \ge\sqrt{\delta_K/32}\,\varepsilon\sqrt N,
\]
contradicting the choice of \(c_K\).
By \eqref{eq:subset-mass}, a fraction at least \(t_0/2\)
of all subsets must therefore have a bad projected block.
This is a pointwise statement for each matrix and needs no
measurable selection of \(v\).
Integrating this finite fraction over the law of \(A\), and using
Proposition~\ref{prop:block}, gives
\[
 \Prob\{s_{\min}(A)<c_K\varepsilon\sqrt N\}
 \le\frac{2m}{k}\Prob_{A,J}
       \{s_{\min}(P_JA_J)<\sqrt{\varepsilon N/8}\}
 \le\frac4{\delta_K\varepsilon}\,q_n
 \longrightarrow0.
\]
The choice of \(c_K\) uses only \(K\), and \eqref{eq:qn}
applies to every admissible order and coordinate-law sequence.
This proves the stated fixed-\(\varepsilon\) convergence.
\end{proof}

\paragraph{AI-use disclosure.}
The argument and manuscript were developed with assistance from
OpenAI Codex. AI assistance was used for literature checking,
calculation checks, drafting, and manuscript review.
No independent human or formal verification is claimed.

\begin{thebibliography}{9}
\bibitem{Vershynin}
R.~Vershynin,
\emph{Concentration inequalities for random tensors},
Bernoulli \textbf{26} (2020), no.~4, 3139--3162.
\href{https://doi.org/10.3150/20-BEJ1218}{doi:10.3150/20-BEJ1218}.
Updated author manuscript:
\href{https://arxiv.org/abs/1905.00802v5}{arXiv:1905.00802v5}.

\bibitem{Tropp}
J.~A.~Tropp,
\emph{An introduction to matrix concentration inequalities},
Foundations and Trends in Machine Learning \textbf{8} (2015),
nos.~1--2, 1--230.
\href{https://doi.org/10.1561/2200000048}{doi:10.1561/2200000048}.
Theorem~5.1.1 in the
\href{https://tropp.caltech.edu/books/Tro14-Introduction-Matrix-preprint.pdf}
{author preprint revised December 24, 2014}.

\bibitem{Fock}
A.~S.~Bandeira, D.~Kunisky, P.~Nizi\'c-Nikolac, L.~Pesenti, and R.~Wang,
\emph{Matrix concentration and equivalent operators on Fock spaces},
\href{https://arxiv.org/abs/2610.01982v1}{arXiv:2610.01982v1} (2026).

\bibitem{vanHandel}
R.~van Handel,
\emph{On the subgaussian comparison theorem},
\href{https://arxiv.org/abs/2512.18588v2}{arXiv:2512.18588v2} (2026;
first posted 2025).

\bibitem{AIM}
S.~Brunton, M.~J.~Colbrook, M.~V.~de Hoop, G.~Stepaniants,
A.~Townsend, and R.~Ward,
\emph{AIM --- Open Applied Problems}, GitHub repository (2026).
Problem 562, revision
\nolinkurl{c9929805cece044705abaaebe1cebddac1ba11a6}.
\href{https://github.com/MColbrook/AIM/blob/c9929805cece044705abaaebe1cebddac1ba11a6/problems/562-tensor-singular-value.md}
{Frozen problem statement}.
\end{thebibliography}
\end{document}
